\chapter[Organización y responsabilidad]{Organización y responsabilidad del departamento de seguridad radiológica}

Dado que la seguridad radiológica es un sistema de prácticas y criterios que involucran todas las actividades que se realizan con material radiactivo, las responsabilidades deben ser compartidas por todo el personal técnico y supervisor. El Departamento de Seguridad Radiológica es el encargado de vigilar que las indicaciones dadas por este manual se lleven a cabo tal y como se especifican.

\section{Organigrama}

Una vista parcial del organigrama de la empresa, mostrando el lugar que ocupa el Departamento de Seguridad Radiológica, se presenta en la figura~\ref{fig:organigrama-seguridad}.

\begin{figure}[htbp]
  \centering
  \begin{tikzpicture}[
    dept/.style={draw, rectangle, align=center, minimum width=4.7cm, minimum height=0.7cm},
    role/.style={draw, rectangle, align=center, minimum width=3.1cm, minimum height=0.7cm},
    node distance=0.65cm and 0.8cm
  ]
    \node (director) [dept] {DIRECTOR GENERAL};
    \node (jefe) [dept, below=of director] {JEFE DEL DEPTO. DE SEGURIDAD\\RADIOLÓGICA};
    \node (auxiliar) [dept, below=of jefe] {AUXILIAR DEL DEPTO. DE SEGURIDAD\\RADIOLÓGICA};
    \node (tecnico) [role, below left=of auxiliar] {TÉCNICO\\RADIÓLOGO};
    \node (auxrad) [role, below right=of auxiliar] {AUXILIAR DEL\\RADIÓLOGO};
    \draw (director) -- (jefe) -- (auxiliar);
    \draw (auxiliar) -- (tecnico);
    \draw (auxiliar) -- (auxrad);
  \end{tikzpicture}
  \caption{Vista parcial del organigrama del Departamento de Seguridad Radiológica.}
  \label{fig:organigrama-seguridad}
\end{figure}

\section{Encargado de seguridad radiológica}

Sus funciones se especifican a continuación:
\begin{enumerate}[label=\roman*)]
  \item Establecer los procedimientos de seguridad radiológica y física aplicables a la adquisición, importación, exportación, producción, posesión, uso, transferencia, transporte, almacenamiento y destino o disposición final de los materiales radiactivos generadores de radiación ionizante, para revisión y aprobación, en su caso, de la CNSNS.

  \item Adiestrar y calificar al POE en la aplicación correcta de las normas y procedimientos de seguridad radiológica y física, así como vigilar su cumplimiento durante las operaciones que se realicen con las fuentes de radiación ionizante.

  \item Establecer el programa de vigilancia radiológica para la determinación, registro, análisis y evaluación de las dosis recibidas por el POE.

  \item Vigilar que al POE se le proporcionen el vestuario, equipo, accesorios y dispositivos de protección radiológica apropiados, y asegurarse de que los use adecuadamente.

  \item Notificar de inmediato a la CNSNS cualquier robo o extravío de las fuentes.

  \item Desarrollar proyectos, procedimientos y métodos para mantener la exposición a la radiación del POE y del público tan baja como razonablemente sea posible, sin exceder los límites especificados en este manual.

  \item Elaborar y supervisar el programa de pruebas de buen funcionamiento y calibración de todo el equipo detector de radiación.

  \item Elaborar, supervisar y participar en los programas de entrenamiento del POE.

  \item Llevar registro de las dosis recibidas por el POE, anexando constancias de dosis recibidas en empleos anteriores cuando las hubiere.

  \item Vigilar que el manejo y la eliminación de los desechos radiactivos se realicen conforme a las normas de seguridad radiológica aplicables.

  \item Efectuar pruebas de fuga a las fuentes de radiación ionizante al momento de su recepción y semestralmente, y llevar un registro de estas pruebas.

  \item Estar presente durante el desarrollo de las diligencias que practique la CNSNS y proporcionar toda la información que le sea requerida.

  \item Expedir al POE los certificados anuales y constancias, al término de la relación laboral, de las dosis recibidas en las 52 semanas anteriores y de la dosis total acumulada a la fecha. Copias de estos documentos, debidamente firmadas de recibido, se enviarán a la CNSNS.

  \item Llevar registro de los exámenes médicos practicados al POE.

  \item Llevar registro de la disposición final de las fuentes agotadas.

  \item Vigilar que en los locales de almacenamiento de material radiactivo se lleven registros de entrada y salida de material radiactivo.

  \item Verificar periódicamente el inventario de material radiactivo.

  \item Otorgar las facilidades que se requieran durante las inspecciones, auditorías, verificaciones y reconocimientos que practique la CNSNS.

  \item Proporcionar la información requerida durante las diligencias anteriores.

  \item Efectuar las pruebas y operaciones que se requieran durante la inspección, auditoría, verificación o reconocimiento.

  \item Corregir las deficiencias y anomalías detectadas en las operaciones anteriores.

  \item Firmar y rubricar toda la documentación que se remita o presente a la CNSNS.

  \item Cubrir todos los gastos derivados de los accidentes radiológicos, incluyendo la indemnización a terceros.

  \item Avisar de inmediato a la CNSNS cuando deje de usar o poseer definitivamente el material radiactivo autorizado.

  \item Tomar todas las medidas de seguridad radiológica y física que se requieran para salvaguardar la integridad de las fuentes de radiación en caso de huelga o paro.

  \item Avisar de inmediato a la CNSNS del estallamiento y terminación de la huelga o paro.
\end{enumerate}

\section{Personal y organización de la CNSNS}

La Comisión Nacional de Seguridad Nuclear y Salvaguardias es la autoridad que reglamenta las actividades desarrolladas para la utilización de la energía nuclear con fines pacíficos.

El personal de la Comisión está organizado en:
\begin{itemize}
  \item \textbf{Gerencia de Seguridad Nuclear.}
  \begin{itemize}
    \item Departamento de Evaluación.
    \item Departamento de Verificación Operativa.
    \item Departamento de Seguridad Física y Salvaguardias.
  \end{itemize}
  \item \textbf{Gerencia de Tecnología, Reglamentación y Servicios.}
  \begin{itemize}
    \item Departamento de Reglamentos, Documentación y Capacitación.
    \item Departamento de Tecnología e Informática.
  \end{itemize}
  \item \textbf{Gerencia de Seguridad Radiológica.}
  \begin{itemize}
    \item Departamento de Supervisión Operativa.
    \item Departamento de Vigilancia Radiológica Ambiental.
  \end{itemize}
  \item Comité de Emergencia.
  \item Consultoría Jurídica.
  \item Oficina de Asuntos Internacionales.
  \item Unidad de Finanzas y Administración.
\end{itemize}
